\documentclass[12pt]{article}

% Layout.
\usepackage[top=1in, bottom=0.75in, left=0.75in, right=0.75in, headheight=1in, headsep=6pt]{geometry}

% Fonts.
\usepackage{mathptmx}

\usepackage[scaled=0.86]{helvet}
\renewcommand{\emph}[1]{\textsf{\textbf{#1}}}

% Misc packages.
\usepackage{amsmath,amssymb,latexsym,amsthm}
\usepackage{graphicx,hyperref}
\usepackage{array}
\usepackage{xcolor}
\usepackage{multicol}
\usepackage{tabularx,colortbl}
\usepackage{enumitem}
\usepackage{soul,multicol}
\usepackage{setspace}
\doublespacing

\hypersetup{
    colorlinks=true,
    linkcolor=blue,
    filecolor=magenta,      
    urlcolor=blue,
    pdftitle={Proofs Worksheet},
    pdfpagemode=FullScreen,
    }

\def\mailto#1{\href{mailto:#1}{#1}}

%% special math symbols
\newcommand{\N}{\mathbb{N}}
\newcommand{\Z}{\mathbb{Z}}
\newcommand{\R}{\mathbb{R}}
\newcommand{\Q}{\mathbb{Q}}
\newcommand{\sse}{\subseteq}
\newcommand{\mcp}{\mathcal{P}}
\newcommand{\ol}[1]{\overline{#1}}

%other specials
\newcommand{\be}{\begin{enumerate}}
\newcommand{\ee}{\end{enumerate}}
\newcommand{\se}{\subseteq}
\newcommand{\spe}{\supseteq}
\newcommand{\ds}{\displaystyle}
\newcommand{\lcm}{\text{lcm}}
%\newcommand{\gcd}{\text{gcd}}

% Paragraph spacing
\parindent 0pt
\parskip 6pt plus 1pt
\def\tableindent{\hskip 0.5 in}
\def\ts{\hskip 1.5 em}

%header
\usepackage{fancyhdr}
\pagestyle{fancy} 
\lhead{\large\sf\textbf{MATH 265: Introduction to Mathematical Proofs}}
\rhead{\large\sf\textbf{Worksheet 16: Proof Critiques}}

\newcommand{\localhead}[1]{\par\smallskip\textbf{#1}\nobreak\\}%
\def\heading#1{\localhead{\large\emph{#1}}}
\def\subheading#1{\localhead{\emph{#1}}}

\newenvironment{clist}%
{\bgroup\parskip 0pt\begin{list}{$\bullet$}{\partopsep 4pt\topsep 0pt\itemsep -2pt}}%
{\end{list}\egroup}%

\begin{document}
Below are four proofs of the assertion: For every pair of sets $A$ and $B$, if $A \subseteq B,$ then $A-B = \emptyset.$\\

\vfill

\textbf{Proof 1:}\\
Suppose $A$ and $B$ are sets and $A \subseteq B.$ So, every element of $A$ is an element of $B$. So, no element can be in $A-B.$ So, $A-B = \emptyset.$\\

\vfill

\textbf{Proof 2:}\\
Suppose $A \subseteq B$. Recall that $A - B$ is the set of all elements that are in $A$ but not in $B$. Since every element of $A$ is also in $B$, there are no elements that are in $A$ but not in $B$.

\medskip
Now suppose for contradiction that $A - B \neq \emptyset$. Then there exists some element $x \in A - B$, which means $x \in A$ and $x \notin B$. But since $A \subseteq B$, we have $x \in B$.

\medskip
Therefore $A - B = \emptyset$.\\

\vfill

\textbf{Proof 3:}\\
(by contrapositive) Suppose $A-B \not= \emptyset.$ Thus, there exists some $x \in A-B.$ Since $x \in A-B,$ it follows that $x \in A$ and $x \not \in B.$ Thus, it is not the case that every element of $A$ is automatically an element of $B.$ Thus, $A \not \subseteq B.$\\

\textbf{Proof 4:}\\
(by contrapositive) Suppose $A-B \not= \emptyset.$ Thus, there exists some $x \in A-B.$ Thus, $x \in A$ and $x \not \in B.$ Thus, $A \not \subseteq B.$\\
\end{document}

\vfill
\item For every pair of sets $A$ and $B,$  $A \subseteq B$ if and only if $A-B = \emptyset.$ 

\begin{proof} Let $A$ and $B$ be sets. We will prove the equivalent statement that $A \not \subseteq B$ if and only if $A-B \not = \emptyset.$ (This is effectively a ``double" contrapositive.)


$( \Rightarrow :)$ Suppose $A \not \subseteq B.$ Then there exists $a \in A$ such that $a \not \in B.$ Thus, $a \in A-B.$ Thus, $A-B \not \emptyset.$

$( \Leftarrow :)$ Suppose $A-B \not= \emptyset.$ Thus, there exists some $x \in A-B.$ Since $x \in A-B,$ it follows that $x \in A$ and $x \not \in B.$ Thus, it is not the case that every element of $x$ is an element of $B.$ Thus, $A \not \subseteq B.$

\end{proof}

\end{enumerate}
\end{document}  















